% Supplementary: exact semantics of TTS + CAC on ReKV (LLaVA-OneVision-7B).
% Requires \usepackage{booktabs} (included in the CVPR template).

\section{Implementation and Semantics of Cross-Frame Token Packing}
\label{sec:supp_semantics}

This section specifies exactly what the language model (LM) receives and attends to
when TTS prunes tokens and CAC packs the survivors into blocks. Unless noted, every
statement applies equally to the unpruned ReKV baseline, with all $P$ tokens of a frame
retained. Hyperparameters are listed in \cref{tab:supp_hparams}.

\paragraph{Frame sampling and timestamps.}
Frames are sampled at $f{=}1$\,fps; frame $t$ is the source frame shown at time $t/f$.
Frames are processed strictly one at a time in arrival order, and no stage reads a
future frame. No explicit timestamp, frame index or separator token is given to the LM:
time is encoded only by the order of tokens in the stream. A question with timestamp
$q$ is answered after frames $0,\dots,\lfloor qf\rfloor$ have been ingested.

\paragraph{Visual tokens and processing order.}
Each frame passes through SigLIP, the multimodal projector, and LLaVA-OneVision's
bilinear $2{\times}$ pooling ($27{\times}27\!\to\!14{\times}14$), giving
$\mathbf{F}_t\in\mathbb{R}^{P\times D}$ with $P{=}196$ tokens in row-major spatial
order. TTS operates \emph{after} pooling, on these LM-space features; the vision tower
and pooling are unchanged.

\paragraph{Token selection (TTS).}
For each spatial position $p$, TTS keeps a reference $\mathbf{r}_p$, the feature last
\emph{retained} at $p$. Token $(t,p)$ is retained iff
$1-\cos(\mathbf{F}_t[p],\mathbf{r}_p)>\tau$, in which case $\mathbf{r}_p\leftarrow
\mathbf{F}_t[p]$; otherwise it is dropped and $\mathbf{r}_p$ is unchanged. Distances are
computed in fp32. The first frame of each video is retained in full and initializes
$\mathbf{r}$; references are reset between videos. Because each decision depends only
on $\mathbf{F}_t$ and $\mathbf{r}$, the selection is causal and independent of how
frames are batched. A dropped token leaves no slot in the sequence: it receives no
attention, and its content survives only in the earlier retained token within cosine
distance $\tau$ at the same position. A frame may retain $0$ to $P$ tokens.

\paragraph{Packing order.}
Retained tokens are appended in $(t,p)$ order: by frame, then by spatial position
within the frame. Tokens from consecutive frames are contiguous, with no padding,
separator or newline token between frames. This is the change in semantics relative to
the unpruned model, where each frame occupies exactly $P$ consecutive positions.

\paragraph{Block formation (CAC).}
The stream of retained tokens is cut into blocks of exactly $B{=}P{=}196$ tokens by
count, regardless of frame boundaries. A block therefore spans a content-dependent
number of frames (many in a static scene, one in a busy one), and a frame spans at most
two blocks. Until $B$ tokens accumulate, retained tokens are buffered as embeddings,
not yet passed through the LM, so they are not yet in the KV cache. Each completed
block is prefilled through the LM and committed to the cache.

\paragraph{Query-time commit.}
When a question arrives, any buffered tokens ($n<B$) are padded to a full block by
repeating the last retained token $B-n$ times. The block is prefilled and committed
before the question is processed, so every frame ingested so far is in memory. The
padded block is permanent: later frames start a new block, and the padding counts
toward the cache size and retrieval cost. On StreamBench this adds on average 96
($\tau{=}0.5$) to 99 ($\tau{=}0.8$) padding tokens per question. At the end of a video
the same padding is applied once to the final partial block.

\paragraph{Warm-up and attention sink.}
Each video starts from an empty cache. The system prompt (13 tokens) is prefilled first
and kept permanently as an attention sink of $n_{\text{init}}$ tokens. While fewer than
$n_{\text{local}}{=}15{,}000$ video tokens have been committed, all of them stay on the
GPU and every block is attended directly. After that, blocks leaving the local window
are offloaded to CPU memory. Pruning only changes \emph{when} this happens: the window
holds $n_{\text{local}}/P\approx 76$ frames unpruned and about
$n_{\text{local}}/(\rho P)$ frames at token keep rate $\rho$.

\paragraph{Position IDs.}
Keys are stored without rotary embedding (RoPE); RoPE is applied at attention time with
positions equal to the index in the attended sequence. During encoding, each new block
attends to the preceding $n_{\text{local}}$ committed tokens at consecutive positions.
Positional distance therefore counts \emph{retained tokens}, not elapsed time: a static
interval that retains few tokens occupies few positions. Following InfLLM/ReKV, once the
sequence exceeds $n_{\text{local}}$ the sink tokens are attended with a single fixed
relative position.

\paragraph{Retrieval.}
Each committed block is represented by the mean of its $B$ keys (padding included);
each layer stores these representatives independently. For a question, the LM first
runs the question tokens alone; in each layer the mean question query is compared with
all block representatives by cosine similarity, and the top $k{=}64$ blocks are
selected per layer. If no blocks have been offloaded yet, the representatives are
computed from the blocks inside the local window and scored by dot product, and all
blocks are used when there are at most $k$. The selected blocks are loaded, sorted by block index (i.e.\ temporal
order) and concatenated after the sink tokens. The prompt and the answer then attend
to [sink; $k$ retrieved blocks; prompt] at consecutive positions, so retrieved blocks
are positionally adjacent regardless of the time between them, exactly as in ReKV. The
retrieved context is always $kB{=}12{,}544$ tokens, but with pruning it covers a
longer, content-dependent stretch of video. The most recent block is eligible for
retrieval like any other, but it is not always selected. Answers are decoded greedily
(at most 256 tokens on StreamBench).

\paragraph{Costs.}
LM prefill during ingestion, KV-cache memory, and the number of retrievable blocks all
scale with the number of committed tokens: retained tokens plus padding, i.e.\ about
$\rho$ of the unpruned count plus about one padding half-block per question. TTS adds
one $D$-dimensional cosine distance per token, negligible beside the LM prefill. The
query-time commit adds at most one block of prefill per question. Retrieval and
answering cost per question are the same as in ReKV, since the retrieved context is a
fixed $kB$ tokens.

\begin{table}[t]
  \centering
  \small
  \begin{tabular}{@{}lr@{}}
    \toprule
    Setting & Value \\
    \midrule
    Backbone & LLaVA-OneVision-Qwen2-7B (fp16) \\
    Sampling rate $f$ & 1\,fps \\
    Tokens per frame $P$ (after $2{\times}$ pooling) & 196 \\
    Block size $B$ & 196 \\
    Local window $n_{\text{local}}$ & 15{,}000 tokens \\
    Attention sink $n_{\text{init}}$ & 13 tokens (system prompt) \\
    Retrieved blocks per layer $k$ & 64 \\
    Block representative & mean key \\
    Retrieval score & cosine (mean question query) \\
    GPU block cache & 128 blocks \\
    TTS distance & $1-\cos$, fp32, vs.\ last retained \\
    Thresholds $\tau$ & 0.50, 0.70 \\
    Decoding & greedy, $\le$256 tokens \\
    \bottomrule
  \end{tabular}
  \caption{Hyperparameters for TTS+CAC on ReKV. All except $\tau$ are identical for the
  unpruned baseline.}
  \label{tab:supp_hparams}
\end{table}
